\subsection{Krull 整环在其分式域有限扩张中的整闭包}

命题 12. 设 A 是 Krull 整环，K 是其分式域，K' 是 K 的有限扩张，A' 是 A 在 K' 中的整闭包。则 A' 是 Krull 整环。A' 的本性赋值是 K' 上与 A 的本性赋值的扩张等价的赋范离散赋值。

设扩张族 \((v_{\iota})_{\iota\in I}\)（扩张到 \(K'\)）由 A 的本性赋值得到。由于次数 \(n = [K':K]\) 有限，赋值 \(v_{\iota}\) 是 \(K'\) 上的离散赋值（第 VI 章，§ 8，第 1 节，命题 1 的推论 3）。令 \(B_{\iota}\) 为 \(v_{\iota}\) 的赋值环；则 \(A' \subset \bigcap_{\iota\in I}B_{\iota}\)（第 VI 章，§ 1，第 3 节，定理 3）。反之，每个元素 \(x\)（属于 \(\bigcap_{\iota\in I}B_{\iota}\)）都整依赖于 A 的每个本性赋值环（第 VI 章，§ 1，第 3 节，定理 3 的推论 3）；因此 \(x\) 在 K 上的极小多项式的系数属于 A（第 V 章，§ 1，第 3 节，命题 11 的推论），所以 \(x\in A'\)；从而 \(A' = \bigcap_{\iota\in I}B_{\iota}\)。现在令 \(x\) 为 \(A'\) 的非零元素；它满足形如 \(x^{s} + a_{s-1}x^{s-1} + \cdots + a_{0} = 0\) 的方程，其中 \(a_{i}\in A\) 且 \(a_{0}\neq 0\)；若 \(v_{\iota}(x) > 0\)，则 \(v_{\iota}(a_{0}) > 0\)；现在，A 的本性赋值 \(v\) 中满足 \(v(a_{0}) > 0\) 的只有有限个，而 \(K'\) 上延拓 K 上给定赋值的赋值也只有有限个（第 VI 章，§ 8，第 3 节，定理 1）；因而 \(v_{\iota}(x) = 0\)，除了有限个指标 \(\iota\in I\) 外。这样便证明了 \(A'\) 是 Krull 整环（第 3 节，定义 3）。

还需证明 \(v_{\iota}\) 与 \(A'\) 的本性赋值等价（第 4 节，命题 6 的推论 2），也就是（第 6 节，定理 3 的推论 2）证明素理想 \(\mathfrak{p}_{\iota}\)（由 \(x \in A'\) 中满足 \(v_{\iota}(x) > 0\) 的元素组成）具有高度 1。若非如此，则存在素理想 \(\mathfrak{q}\)，它是 \(A'\) 的素理想，并满足 \((0) \subset \mathfrak{q} \subset \mathfrak{p}_{\iota}\)，且不同于 (0) 和 \(\mathfrak{p}_{\iota}\)；于是 \((0) \subset \mathfrak{q} \cap A \subset \mathfrak{p}_{\iota} \cap A\)，并且 \(\mathfrak{q} \cap A\) 不同于 (0) 和 \(\mathfrak{p}_{\iota} \cap A\)（第 V 章，§ 2，第 1 节，命题 1 的推论 1）；因此素理想 \(\mathfrak{p}_{\iota} \cap A\) 的高度不为 1，这与它对应于 \(A\) 的本性赋值这一事实矛盾。

推论。设 \(\mathfrak{p}\)（分别地，\(\mathfrak{p}'\)）是 A（分别地，A'）的高度为 1 的素理想，\(v\)（分别地，\(v'\)）是与之对应的 A（分别地，A'）的本性赋值。\(\mathfrak{p}'\) 位于 \(\mathfrak{p}\) 之上的充要条件是 \(v'\) 在 K 上的限制等价于 \(v\)。

赋值 \(v'\) 等价于 A 的本性赋值 w 的扩张（命题 12）。令 \(q = p' \cap A\)，它是 A 的高度为 1 的素理想。\(v'\) 在 K 上的限制等价于 v 的充要条件是 w = v，从而 q = p。

